% ============================================================
\section{Drivers y su priorización}\label{sec:priorizacion}
% ============================================================

Esta sección deriva los drivers de los ASR de la sección~\ref{sec:asr} y los ordena. Un driver agrupa los ASR que dependen de la misma decisión estructural; su prioridad empieza por la de sus ASR y se completa con el impacto de esa decisión sobre la arquitectura. Donde el documento de contexto y la sección~\ref{sec:asr} difieren, manda la sección~\ref{sec:asr}.

% ------------------------------------------------------------
\subsection{De los ASR a los drivers}

\textbf{Criterio de agrupación.} Dos ASR van al mismo driver si una sola decisión estructural los cumple o los incumple a la vez. Un ASR puede alimentar a más de un driver cuando exige dos decisiones distintas: ASR-04 necesita que el estado sobreviva a la conexión y además un protocolo que permita reasociarla, y ASR-14 fija a la vez el protocolo, la persistencia y el lenguaje del motor.

\begin{tabla}{|L{1.4cm}|Y|L{3.3cm}|L{3.6cm}|}
\hline
\filacab \cab{Driver} & \cab{Enunciado} & \cab{ASR que agrupa} & \cab{Decisión estructural común} \\ \hline
\endhead
DRV-01 & Servidor autoritativo: el servidor arbitra cada partida en línea y es la única autoridad sobre su estado; el cliente propone y presenta; el servidor decide y conserva el estado vigente, con independencia de las conexiones y de la base de datos. & ASR-08, ASR-02, ASR-04, ASR-05 (prioritarios); ASR-09 (media) & Dónde vive el estado vigente de una partida y quién lo produce. \\ \hline
DRV-02 & Protocolo propio sobre TCP: sin WebSockets, el equipo diseña un protocolo que resuelve la delimitación, el versionado, la validación, la sesión y el cifrado, con un costo por mensaje compatible con el segundo por jugada. & ASR-03, ASR-04, ASR-14 (prioritarios) & Cómo viaja un mensaje entre cliente y servidor. \\ \hline
DRV-03 & Presupuesto de un segundo por jugada: menos de un segundo, con mediana bajo 300 ms, repartido de forma explícita entre los tramos del camino de la jugada. & ASR-01 (prioritario); ASR-09 (media) & Cuánto tiempo consume cada tramo del camino de una jugada. \\ \hline
DRV-04 & Política explícita de persistencia por tipo de dato: qué se escribe de forma confirmada, qué de forma diferida con ventana acotada y qué constancia queda de lo perdido, con la capa de datos aislada. & ASR-05, ASR-14 (prioritarios) & Qué se escribe, cuándo y qué se pierde ante una caída. \\ \hline
DRV-05 & Un único motor de reglas, parametrizable, determinista y ejecutable sin interfaz. & ASR-06, ASR-14 (prioritarios) & Dónde viven las reglas y sus parámetros. \\ \hline
DRV-06 & Protección del menor: control en el servidor, acceso que un niño puede completar y datos mínimos. & ASR-11, ASR-12 (alta) & Qué controles viven en el servidor y qué se guarda del menor. \\ \hline
DRV-07 & Localización de extremo a extremo: todo texto y formato se adapta sin tocar código, en Unicode de la pantalla a la base de datos. & ASR-13 (alta) & Dónde viven los textos y los formatos. \\ \hline
DRV-08 & Plazo de 14 semanas con un equipo de dos personas: rebanadas verticales, regresión semanal sin interfaz y puntos de variación acotados a lo que puede recortarse. & ASR-15 (alta); ASR-10 (media); ASR-07 (reservado) & Qué se aísla para poder recortarse y en qué orden se construye. \\ \hline
\end{tabla}

\textbf{Qué cambia respecto al documento de contexto.} Resultan los mismos ocho drivers, pero sus enunciados se ajustan a los ASR. DRV-01 incorpora la partida en línea arbitrada (ASR-08), que el documento de contexto daba por supuesta. DRV-06 y DRV-07 dejan de ser drivers sin ASR: ahora recogen ASR que proceden de restricciones y de funcionalidad. DRV-08 incorpora el tamaño del equipo (ASR-15) y la IA como funcionalidad recortable (ASR-10). ASR-01 conserva su medida, pero la decisión que forzaba en el documento de contexto, validar parte de la jugada en el cliente, queda descartada por DEC-02.

\textbf{Cobertura.} Los catorce ASR activos quedan dentro de algún driver, y el reservado ASR-07 en DRV-08. Ningún driver existe sin al menos un ASR que lo sostenga.

% ------------------------------------------------------------
\subsection{Criterios}

Cada driver se califica con cuatro criterios, de 1 a 5. El primero viene directamente de la sección~\ref{sec:asr}; los otros tres miden el impacto de la decisión que el driver agrupa. El total ordena las iteraciones.

\begin{tabla}{|L{3.4cm}|Y|}
\hline
\filacab \cab{Criterio} & \cab{Qué mide} \\ \hline
\endhead
Peso de sus ASR &
La prioridad de los ASR que agrupa, tal como la fijó la subsección~\ref{sec:asr:prioridad}. Un 5 si agrupa tres o más ASR prioritarios; un 4 si agrupa dos; un 3 si agrupa uno; un 2 si solo agrupa ASR de prioridad alta; y un 1 si solo agrupa ASR de prioridad media. \\ \hline
Alcance estructural &
Cuántos elementos del sistema cambian si el driver se resuelve de otra manera. Un 5 significa que cambian a la vez cliente, servidor, protocolo y capa de datos; un 1, que el efecto queda dentro de un componente. \\ \hline
Costo de cambio tardío &
Qué cuesta cambiar la decisión cuando ya hay código construido. Pesa aquí que haya clientes desplegados, datos escritos con un formato o componentes construidos sobre el supuesto anterior. \\ \hline
Drivers que habilita &
Cuántos de los demás drivers dependen de que este ya esté resuelto para poder decidirse. Un driver que no habilita a ninguno puede esperar sin bloquear el avance. \\ \hline
\end{tabla}

% ------------------------------------------------------------
\subsection{Calificación y orden de las iteraciones}

\begin{tabla}{|L{1.4cm}|L{3.6cm}|L{1.3cm}|L{1.3cm}|L{1.3cm}|L{1.3cm}|L{1.0cm}|L{1.2cm}|}
\hline
\filacab \cab{Driver} & \cab{Enunciado abreviado} & \cab{Peso ASR} & \cab{Alcance} & \cab{Costo tardío} & \cab{Habilita} & \cab{Total} & \cab{Iter.} \\ \hline
\endhead
DRV-01 & Servidor autoritativo & 5 & 5 & 5 & 5 & 20 & 1 \\ \hline
DRV-02 & Protocolo propio sobre TCP & 5 & 4 & 5 & 4 & 18 & 2 \\ \hline
DRV-05 & Motor de reglas único & 4 & 4 & 4 & 4 & 16 & 3 \\ \hline
DRV-03 & Presupuesto de un segundo & 3 & 5 & 4 & 3 & 15 & 4 \\ \hline
DRV-04 & Política de persistencia & 4 & 3 & 4 & 3 & 14 & 5 \\ \hline
DRV-06 & Protección del menor & 2 & 3 & 3 & 2 & 10 & 6 \\ \hline
DRV-07 & Localización de extremo a extremo & 2 & 3 & 4 & 1 & 10 & 7 \\ \hline
DRV-08 & Plazo, equipo y puntos de variación & 2 & 2 & 2 & 2 & 8 & 8 \\ \hline
\end{tabla}

\textbf{Cómo se rompe el empate.} DRV-06 y DRV-07 empatan en 10 y se ordenan por el peso de sus ASR: DRV-06 agrupa dos ASR de prioridad alta y DRV-07 uno. Además, ASR-11 y ASR-12 nacen de requisitos impuestos con stakeholders de veto (STK-02, STK-17), mientras que ASR-13 nace de un objetivo de negocio.

\textbf{El orden sigue a los ASR.} Los cinco drivers que agrupan ASR prioritarios ocupan las cinco primeras iteraciones, y los tres que solo agrupan ASR de prioridad alta o media, las tres últimas. Dentro de cada grupo decide el impacto: DRV-03 tiene más alcance que DRV-05, pero va después porque no se puede repartir el segundo hasta que el trabajo del motor esté acotado, y eso ya lo refleja su menor calificación en drivers habilitados.

\textbf{DRV-08 se aplica además de forma transversal.} Es el último en tener iteración propia porque no cambia qué componentes existen, pero su exigencia de rebanadas verticales gobierna el orden en que se construye lo que las demás iteraciones deciden. Su iteración se ocupa de acotar los puntos de variación, no de descubrirlos.

% ------------------------------------------------------------
\subsection{Plan de iteraciones}\label{sec:priorizacion:plan}

Este documento desarrolla la primera iteración de cada driver, en el orden de la tabla anterior. Cada una satisface un ASR concreto de la sección~\ref{sec:asr}, con un solo objetivo, un solo elemento y una sola medida, y termina en un incremento demostrable.

\begin{tabla}{|L{0.9cm}|L{1.4cm}|L{1.4cm}|L{4.4cm}|Y|}
\hline
\filacab \cab{Iter.} & \cab{Driver} & \cab{ASR} & \cab{Objetivo único} & \cab{Elemento a descomponer y medida} \\ \hline
\endhead
1 & DRV-01 & ASR-02 & El servidor es la única autoridad sobre el estado de una partida en curso. & Gestión del estado de la partida. Medida: ESC-05. \\ \hline
2 & DRV-02 & ASR-03 & El tráfico capturado no revela nada legible ni se puede reutilizar. & Servicio de protocolo. Medida: ESC-06. \\ \hline
3 & DRV-05 & ASR-06 & Las reglas cambian sin tocar código y el motor se ejecuta sin interfaz. & Motor de reglas. Medida: ESC-12. \\ \hline
4 & DRV-03 & ASR-01 & Una jugada se responde en menos de un segundo. & Camino de la jugada. Medida: ESC-01. \\ \hline
5 & DRV-04 & ASR-05 & Una caída de la base de datos pierde solo lo acordado y deja constancia. & Servicio de acceso a datos. Medida: ESC-09. \\ \hline
6 & DRV-06 & ASR-11 & Ningún mensaje inapropiado llega a un menor, ni siquiera cuando el filtro falla. & Servicio de chat. Medida: ESC-04. \\ \hline
7 & DRV-07 & ASR-13 & Se agrega un idioma nuevo sin tocar código ni rehacer pantallas. & Catálogo de textos y formatos. Medida: ESC-11. \\ \hline
8 & DRV-08 & ASR-15 & El oponente controlado por IA se puede recortar sin tocar el núcleo. & Punto de variación del oponente. Medida: ESC-13. \\ \hline
\end{tabla}

\subsubsection*{Cómo quedan atendidos los ASR}

\begin{tabla}{|L{1.4cm}|L{3.2cm}|Y|}
\hline
\filacab \cab{ASR} & \cab{Iteración} & \cab{Papel que cumple} \\ \hline
\endhead
ASR-02 & 1 & Objetivo. \\ \hline
ASR-08 & 1 & Base del incremento: la partida en línea arbitrada que CU-20 y CU-21 recorren de principio a fin. \\ \hline
ASR-03 & 2 & Objetivo. \\ \hline
ASR-14 & 2, 3 y 5 & Condición: el protocolo propio en la~2, el motor en C\# en la~3 y SQL Server en la~5. \\ \hline
ASR-06 & 3 & Objetivo. \\ \hline
ASR-01 & 4 & Objetivo. \\ \hline
ASR-09 & 4 & Condición: la notificación a los espectadores no puede frenar a los jugadores. \\ \hline
ASR-05 & 5 & Objetivo. \\ \hline
ASR-11 & 6 & Objetivo. \\ \hline
ASR-13 & 7 & Objetivo. \\ \hline
ASR-15 & 8 & Objetivo. \\ \hline
ASR-10 & 8 & Condición: la IA existe como un participante que se puede retirar. \\ \hline
ASR-04 & Pendiente & Necesita la sesión de partida, que no se puede decidir antes del protocolo de la iteración~2. La iteración~2 le deja sitio en el apretón de manos. \\ \hline
ASR-12 & Pendiente & Necesita el módulo de cuentas; es la segunda iteración de DRV-06. \\ \hline
ASR-07 & Reservado & Solo si Q-02 se resuelve a favor. \\ \hline
\end{tabla}

\textbf{Lo que falta después de estas ocho.} Un driver con varios ASR necesita más de una iteración, y aquí solo está la primera de cada uno. Quedan dos ASR sin iteración propia: ASR-04, que comparten DRV-01 y DRV-02, y ASR-12, de DRV-06. Quedan también los otros dos puntos de variación de DRV-08, el extremo de red, que se volverá obligatorio si ASR-07 se activa, y la marca con los recursos de terceros.

% ------------------------------------------------------------
\subsection{Por qué DRV-01 va primero}

DRV-01 obtiene la calificación máxima en los cuatro criterios, y es el único que la obtiene.

\begin{itemize}
  \item \textbf{Peso de sus ASR.} Agrupa cuatro ASR prioritarios, entre ellos los dos que abren el orden de la sección~\ref{sec:asr}: ASR-08, sin el cual no hay servidor, y ASR-02, que decide que el servidor tiene la autoridad.
  \item \textbf{Alcance estructural.} Decide quién produce el estado de la partida y quién lo consume. Cambiar de servidor autoritativo a cliente autoritativo obliga a rehacer a la vez el cliente, el servidor, el protocolo y la capa de datos: cada mensaje cambia de significado, porque deja de ser una propuesta y pasa a ser un hecho.
  \item \textbf{Costo de cambio tardío.} Es la decisión que CRN-10 describe como la más cara de cambiar, y ASR-02, ASR-04 y ASR-05 fallan a la vez si el estado vigente no tiene un dueño único en el servidor.
  \item \textbf{Drivers que habilita.} DRV-02 no puede fijar el apretón de manos sin saber que el estado sobrevive a la conexión; DRV-03 no puede repartir el segundo sin saber cuántos viajes hay por jugada; DRV-04 no puede definir qué se escribe y cuándo sin saber quién tiene el estado vigente; y DRV-06 aplica el filtro de chat donde vive la autoridad.
\end{itemize}
